% Method: formal mapping objects used by the scoped P3 result, plus clearly marked
% broader framework objects that remain untested externally.
This section defines the scientific mapping objects used by the scoped
public-reference evaluation. Their original broader design lineage is preserved
in \texttt{P3.cross-domain-atlas.v1}; the executed mapping protocol is the
additive \texttt{P3.public-reference-mapping.v1.1-confirmatory}. The definitions
are data-level and do not depend on a particular language model or retrieval
backend.

\subsection{Source and source-local projection}

Let a source record $s$ bind an external document or structured source item to
an exact identity, version/revision, locator, and content commitment. ORION
represents the relevant scientific meaning as a source-local projection
\begin{equation}
    P_s = \langle R_s, C_s, M_s, X_s, \rho_s, \mu_s, a_s, \delta_s, \Pi_s \rangle,
\end{equation}
with coordinates for referent $R_s$, construct $C_s$,
measurement/operationalization $M_s$, context $X_s$, polarity $\rho_s$,
modality $\mu_s$, attribution $a_s$, discourse relation $\delta_s$, and
assumptions/unresolved information $\Pi_s$.

Each coordinate is a typed hypothesis supported by the source authority
available for that case, not a value filled merely to complete a schema. In the
scoped public-reference datasets, unsupported coordinates remain absent or
unresolved; the study does not ask a model to infer missing gold coordinates
from raw papers.

\subsection{Coordinate-wise comparison}

A comparison between projections $P_a$ and $P_b$ is a vector of typed
coordinate relations rather than a scalar similarity score:
\begin{equation}
\begin{split}
    q(P_a,P_b) = \langle &\mathrm{ref}(R_a,R_b),
    \mathrm{const}(C_a,C_b),
    \mathrm{meas}(M_a,M_b),
    \mathrm{ctx}(X_a,X_b),\\
    &\mathrm{pol}(\rho_a,\rho_b),
    \mathrm{mod}(\mu_a,\mu_b),
    \mathrm{attr}(a_a,a_b),
    \mathrm{disc}(\delta_a,\delta_b) \rangle .
\end{split}
\end{equation}
Relations can express compatibility, incompatibility, conditional
transformability, or unresolved evidence according to the frozen semantic
taxonomies. A proposed relation does not acquire authority merely from lexical
or embedding similarity.

\subsection{Typed mapping and preservation conditions}

A typed mapping $m$ states which coordinate relations it licenses and under
what preservation conditions. Its record contains:

\begin{itemize}
    \item the proposed relation (for example identity, equivalence,
          conditional transform, subsumption, association-only,
          no-licensed-mapping, or unresolved);
    \item the coordinates the mapping claims to preserve;
    \item conditions that must hold for the mapping to be used; and
    \item coordinates or information that are not preserved and therefore may
          not be silently transferred into a global statement.
\end{itemize}

This representation is intentionally stricter than a successful match score:
a mapping can be structurally plausible yet inadmissible for the scientific
claim if a load-bearing condition is absent.

\subsection{GLUE, obstruction, and unresolved mapping}

Two or more projections may be \textsc{glue}d only when the coordinate
relations required by the mapping are compatible and its preservation
conditions are satisfied. If a required relation is incompatible, the correct
output is a typed \obstruction{} naming the failing coordinate or condition. If
the available evidence is insufficient, the relation remains \unresolved{}
rather than being forced to merge or split. Multiple locally valid but
mutually incompatible views may remain a \pluralview{}.

The scoped confirmatory experiment tests this decision semantics directly. Its
strongest discrimination is on polarity/modality/attribution/context cases,
where flat predicate canonicalization merges claims that the typed rule keeps
separate. It does not identify the necessity of every coordinate; coordinate
necessity requires targeted cases in follow-up issue \#280.

\subsection{Provenance and recoverable records}

Every scoped case retains the external source identity and the frozen record
used to derive the expected mapping decision. A mapping result retains the
source records and the coordinate conditions supporting its terminal. This
provenance is necessary for independent replay of P3.C5.

The broader ORION design goes further and defines a global portrait
\begin{equation}
    \mathcal{G} = \langle \{P_i\}, \mathcal{M}, \sigma \rangle,
\end{equation}
where $\mathcal{M}$ stores accepted and rejected mappings and $\sigma$ maps a
global statement back to source projections and conditions. Reader-level
recoverability of generated portraits is a follow-up empirical claim (P3.C8/
broader protocol), not a result of the current mapping study.

\subsection{Search-universe feedback is outside the scoped empirical claim}

The wider ORION architecture allows newly absorbed representations to change
the search universe $W$ and reopen future acquisition routes. That feedback is
part of the broader programme but is not manipulated or measured in the
public-reference mapping experiment. No result in this paper is attributed to
a ``W effect,'' and no W-expansion ablation is used to support P3.C5.

\subsection{What is and is not claimed}

ORION does not claim the first scientific knowledge graph, cross-domain RAG
system, schema matcher, provenance graph, entity resolver, or
literature-discovery method. The supported empirical object is narrower: over
the frozen public-reference mapping cases, explicit scientific-coordinate
conditions plus an obstruction state reduce false merges relative to flat
predicate canonicalization while satisfying a conservative false-split guard.

Local exact worlds~\cite{orion2025paper1,orion2025paper2} continue to falsify
implementation semantics in controlled cases. They do not establish external
raw-text extraction quality, recoverability, downstream utility, or broad
cross-domain generality; those remain separate prospective questions.

\endinput
